\section{Test Case 15: forced array reallocation}

Test Case 15 is a dynamic allocation check. It starts with 100 beads and no
spare initial capacity. The nozzle velocity is set to 200,000 cm/s to produce
repeated insertions in the short run. The first insertion exceeds
\texttt{mxnpjet} and calls
\texttt{reallocate\_jet}. The program reports \texttt{Array reallocations:}
at shutdown; the acceptance condition is a value greater than zero and at
least one topology addition. Each subsequent reallocation grows the capacity
by 100 slots; the initial allocation remains governed by the normal
100-slot increment. The small persistent GPU topology path is also enabled so
that capacity growth and device remapping are exercised. Acceptance checks
completion, topology additions, and a positive reallocation count rather than
bead-for-bead trajectory identity. The input is stored in
\texttt{examples/input-15/input.dat}.

The relaxed acceptance criterion was originally attributed to the high
insertion speed amplifying GPU rounding. That explanation was incorrect. The
divergence was caused by a defect in the accelerator path: on a reallocation
the persistent device mapping of the Coulomb force array was not released
before the host array changed, so the device either failed with a
\textit{partially present} error or silently reused the previous buffer. The
GPU run produced NaN from its first printed line with the bead count frozen at
103, against 211 beads and a final abscissa of 35.99 cm on the CPU. Once the
mapping is released correctly, the two agree to eight significant digits. This
test is consequently a genuine regression guard for that defect, and a future
tightening of its acceptance criterion is justified.
